\subsection{Staging và ODS: làm sạch dữ liệu}
Dữ liệu từ hai hệ thống nguồn trước hết được chép nguyên trạng vào tầng staging (11 bảng), mỗi dòng gắn mã đợt nạp và thời điểm nạp để truy vết. Có hai chế độ trích xuất: toàn phần (xoá và chép lại) và tăng dần (chỉ lấy dòng mới dựa trên watermark: ngày mua của đơn, ngày liên hệ đầu của MQL, ngày chốt của deal; các bảng con lấy theo đơn mới). Từ staging, các thủ tục SQL làm sạch dữ liệu vào tầng ODS. Bảng toạ độ địa lý 1.000.163 điểm được gom về 19.010 mã bưu chính, lấy median toạ độ và mode tên thành phố sau khi loại 42 điểm nằm ngoài phạm vi Brazil. Tên thành phố được chuẩn hoá chữ thường, bỏ dấu (73.440 điểm toạ độ bị đổi). Cột sai chính tả \code{*\_\allowbreak{}lenght} được đổi tên; 610 sản phẩm không có danh mục được gán \code{unknown}; 13 sản phẩm thuộc hai danh mục thiếu bản dịch được bổ sung bằng bảng seed. Với đơn hàng, thay vì xoá, hệ thống gắn cờ lỗi: 129 trong 45.430 đơn có mốc thời gian sai thứ tự hoặc thiếu ngày (phổ biến nhất là ngày bàn giao vận chuyển trước ngày duyệt, 84 đơn). Các dòng không nạp được đều ghi vào bảng \code{etl\_\allowbreak{}log.rejected\_\allowbreak{}rows} kèm lý do; kiểm thử đối soát cho thấy với mọi bảng số dòng staging bằng số dòng ODS cộng số dòng bị loại.

Các quy tắc làm sạch và số dòng bị ảnh hưởng được tổng hợp ở Bảng \ref{tab:dq-errors} (mục 3.3.3). Mỗi dòng bị loại được ghi vào \code{etl\_log.rejected\_rows} cùng lý do; không có dòng nào bị bỏ âm thầm.
\todo{Thêm sơ đồ luồng staging $\rightarrow$ ODS (các job ETL) và bảng ``nguồn $\rightarrow$ đích $\rightarrow$ phép biến đổi'' (Phụ lục C).}
